\chapter{Introducción}
<<Aquí va la introducción del TFG. Extensión máxima: 10 páginas.>>\\

\textbf{Contenido de la sección:}
\begin{itemize}
    \item Introducción al tema del TFG y las preguntas respondidas en él
    \item Contextualización
    \begin{itemize}
        \item Contexto general
    \end{itemize}
    \item Motivación
    \begin{itemize}
        \item Problemática detectada a resolver
        \item Relevancia en el contexto actual
    \end{itemize}
    \item Objetivos
    \begin{itemize}
        \item Objetivos generales y específicos
        \item Hipótesis / Pregunta de investigación
    \end{itemize}
    \item (Estructura del manuscrito) <- Solo si es necesario
    \begin{itemize}
        \item Qué hay en cada capítulo
    \end{itemize}
\end{itemize}

\textbf{= OLD =}


Esta plantilla muestra la estructura básica de la memoria final del TFG, así como algunas instrucciones de formato.

El \textbf{esquema básico de una memoria} final \textbf{de TFG} es el siguiente:
\begin{itemize}
\item[•] Resumen en español [máximo 2 páginas]
\item[•] Abstract en inglés [máximo 2 páginas]
\item[•] Tabla de contenidos
\item[•] Introducción [máximo 10 páginas]
\begin{itemize}
    \item Contextualización del proyecto
    \item Motivación para su realización
\end{itemize}
\begin{itemize}
    \item Objetivos del Proyecto
\end{itemize}
\item[•]   Planificación [máximo 5 páginas]
\begin{itemize}
    \item Planificación de tiempos
    \item Costes
\end{itemize}
\begin{itemize}
    \item Casos de Uso
    \item Requisitos del Proyecto
    \item Diseño y Pruebas
\end{itemize}
\item[•] Estado del Arte [máximo 15 páginas]
\begin{itemize}
    \item Análisis del mercado
    \item Análisis de las tecnologías empleadas
\end{itemize}
\item[•] Desarrollo [máximo 60 páginas]
\begin{itemize}
    \item Desarrollo
    \item Resultados preliminares y Verificación (Cumpliendo los requisitos)
\end{itemize}
\item[•] Resultados [máximo 4 páginas]
\begin{itemize}
    \item Validación (Condiciones de uso reales): Resultados finales
\end{itemize}
\item[•] Conclusiones [máximo 1 página]
\begin{itemize}
    \item Conclusiones
    \item Trabajos Futuros
\end{itemize}
\item[•] Objetivos de Desarrollo Sostenible (ODS) y Aspectos Sociales [máximo 1 página]
\begin{itemize}
    \item ODSs
    \item Impacto Social y Medioambiental 
\end{itemize}
\item[•] Bibliografía (publicaciones utilizadas en el estudio y desarrollo del Trabajo)
\item[•] Anexos (opcional)
\end{itemize}

En cualquier caso, es el tutor del TFG quien indicará a su estudiante la estructura final de la memoria que mejor se ajuste al trabajo realizado.

Con respecto al formato, se seguirán las siguientes pautas, que se muestran en esta plantilla:
\begin{itemize}
\item[•] \textit{\textbf{Tamaño de papel}:} DIN A4
\item[•] \textit{\textbf{Portada}:} Tal y como se recoge en esta plantilla, con la indicación de la universidad, centro, título del TFG y autor.
\item[•] \textit{\textbf{Segunda página}:} Información bibliográfica, incluyendo todos los datos del Tutor del TFG y del codirector (si lo hubiera).
\item[•] \textit{\textbf{Tipo de letra para el texto}:} Preferiblemente “Bookman Old Style” de 11 puntos. Si no fuera posible, las alternativas recomendadas son, por orden de preferencia: “Palatino Linotype”, “Garamond” o “Georgia”.
\item[•] \textit{\textbf{Tipo de letra para código fuente}:} “Consolas” o “Roboto mono”
\item[•] \textit{\textbf{Márgenes}:} Superior e inferior $3$ cm, izquierdo y derecho $2.54$ cm.
\item[•] \textit{\textbf{Secciones y subsecciones}:} Reseñadas con numeración decimal a continuación del número del capítulo. Ej.: subsecciones 2.3.1.
\item[•] \textit{\textbf{Números de página}:} Siempre centrados en el margen inferior. La página 1 comienza en el capítulo 1, y todas las secciones anteriores al capítulo 1 se numeran en números romanos en minúscula (i, ii, iii…).
\end{itemize}

\vspace*{1.5cm}
Para elaborar la memoria final del TFG con esta plantilla, seguir los siguientes pasos:
\begin{enumerate}
%\item Descargar e instalar MiKTeX:  \url{https://miktex.org/}
%\item Descargar e instalar un editor de \LaTeX~, por ejemplo Texmaker:\\
%\url{https://www.xm1math.net/texmaker/}

\item Editar el archivo \textbf{secciones/ \_DatosTFM.tex}, que está en la carpeta \textbf{secciones} de esta plantilla. Cumplimentar todos los datos solicitados en dicho archivo. Guardar y cerrar.
\item Compilar el archivo \textbf{plantilla\_TFG.tex} (puede ser renombrado). Se generará como resultado un archivo \textbf{pdf}.
\item Para escribir la memoria final del TFG, se pueden añadir y/o modificar los archivos de la carpeta \textbf{secciones} según sea necesario. El resultado se obtiene al compilar el archivo \textbf{plantilla\_TFG.tex}, por lo que se debe modificar en consecuencia.
\end{enumerate}


\section{Ejemplo de código en Python}
\begin{lstlisting}[style=Python]
import numpy as np
from scipy.spatial.transform import Rotation

def main(args=None):
    print('Hola mundo!')

if __name__ == '__main__':
    main()
\end{lstlisting}

\section{Ejemplo de código en C++}
\begin{lstlisting}[style=C++]
#include <iostream>
#include <vector>
using namespace std;

int main() {
    vector<int> v = {1, 2, 3};
    for (int i = 0; i < 10; i++) {
        cout << i << " ";
    }
    cout << endl;
    return 0;
}
\end{lstlisting}
